\chapter*{Listă de abrevieri și termeni}
\addcontentsline{toc}{chapter}{Listă de abrevieri și termeni}

Tabelele de mai jos rezumă abrevierile tehnice și termenii din domeniul apicol folosiți pe parcursul lucrării.

\section*{Abrevieri și acronime}

\begin{longtable}{@{}p{2.6cm} p{10.6cm}@{}}
\textbf{Abreviere} & \textbf{Semnificație} \\
\hline
\endfirsthead
\textbf{Abreviere} & \textbf{Semnificație} \\
\hline
\endhead
AI & Inteligență artificială: domeniul care urmărește construirea de sisteme capabile să execute sarcini ce necesită în mod normal inteligență umană, precum recunoașterea imaginilor. \\
API & (\textit{Application Programming Interface}) Contractul prin care două programe comunică între ele: setul de funcții, puncte de acces și reguli prin care o aplicație poate folosi serviciile alteia. \\
CNN & (\textit{Convolutional Neural Network}) Rețea neuronală specializată în prelucrarea imaginilor, care învață singură tipare vizuale prin filtre aplicate succesiv asupra pixelilor. \\
CRUD & Acronim pentru cele patru operații fundamentale asupra datelor: creare, citire, actualizare și ștergere (\textit{Create, Read, Update, Delete}). \\
DAO & (\textit{Data Access Object}) Componentă software care concentrează și izolează codul de acces la baza de date față de restul aplicației. \\
DI & (\textit{Dependency Injection}) Mecanism prin care un obiect primește din exterior componentele de care depinde, în loc să le construiască singur. \\
DTO & (\textit{Data Transfer Object}) Obiect simplu, fără logică de business, folosit pentru a transporta date prin rețea sau între straturile unei aplicații. \\
EF Core & (\textit{Entity Framework Core}) Instrumentul Microsoft de mapare obiect-relațională pentru .NET, care permite lucrul cu baza de date prin obiecte C\# și generează automat comenzile SQL. \\
FK & Cheie externă (\textit{Foreign Key}): coloană care leagă o entitate de cheia primară a alteia și materializează relațiile dintre tabele. \\
HTTP / HTTPS & Protocolul de comunicație al aplicațiilor web; varianta HTTPS adaugă un strat de criptare peste HTTP. \\
IoT & (\textit{Internet of Things}) Rețea de dispozitive fizice echipate cu senzori și conectate la internet, care colectează și schimbă date. \\
IP & (\textit{Internet Protocol}) Adresa numerică ce identifică un dispozitiv într-o rețea. \\
JSON & (\textit{JavaScript Object Notation}) Format text, ușor de citit atât de oameni, cât și de programe, folosit pe scară largă pentru schimbul de date. \\
JVM & (\textit{Java Virtual Machine}) Mediul de execuție pe care rulează codul compilat Java și Kotlin. \\
JWT & (\textit{JSON Web Token}) Token semnat criptografic care atestă identitatea unui utilizator autentificat. \\
MVVM & (\textit{Model–View–ViewModel}) Șablon arhitectural care separă afișarea (View) de logica de prezentare (ViewModel) și de date (Model). \\
ORM & (\textit{Object-Relational Mapping}) Tehnică prin care obiectele din cod sunt puse în corespondență cu tabelele unei baze de date relaționale. \\
PDF & (\textit{Portable Document Format}) Format de document independent de platformă, destinat vizualizării și tipăririi consecvente. \\
PK & Cheie primară (\textit{Primary Key}): coloana care identifică unic fiecare rând dintr-un tabel. \\
QR & (\textit{Quick Response}) Cod de bare bidimensional care poate fi citit rapid cu camera unui telefon. \\
ReLU & (\textit{Rectified Linear Unit}) Funcție de activare folosită în rețelele neuronale, care păstrează valorile pozitive și anulează valorile negative. \\
REST & (\textit{Representational State Transfer}) Stil de proiectare a serviciilor web, organizat în jurul resurselor și al metodelor HTTP standard. \\
SDK & (\textit{Software Development Kit}) Ansamblu de instrumente, biblioteci și documentație pentru dezvoltarea aplicațiilor pe o anumită platformă. \\
SMTP & (\textit{Simple Mail Transfer Protocol}) Protocolul standard folosit pentru trimiterea mesajelor de e-mail. \\
SQL & (\textit{Structured Query Language}) Limbajul standard de interogare și manipulare a bazelor de date relaționale. \\
TTL & (\textit{Time To Live}) Durata cât o valoare rămâne validă, de exemplu într-un cache, înainte de a fi considerată expirată. \\
UI & (\textit{User Interface}) Interfața grafică prin care utilizatorul interacționează cu o aplicație. \\
URL / URI & Adresa care identifică și localizează o resursă; un „Data URI" conține resursa (de exemplu, o imagine) codată direct în adresă. \\
UUID & (\textit{Universally Unique Identifier}) Identificator de 128 de biți, practic unic la nivel global, folosit pentru a eticheta entități fără coordonare centrală. \\
\end{longtable}

\section*{Termeni din domeniul apicol}

\begin{longtable}{@{}p{2.6cm} p{10.6cm}@{}}
\textbf{Termen} & \textbf{Explicație} \\
\hline
\endfirsthead
\textbf{Termen} & \textbf{Explicație} \\
\hline
\endhead
Apicultură de precizie & Direcție de gestionare a coloniilor pe baza datelor colectate, pentru decizii mai bune și pierderi reduse. \\
Botcă & Celulă specială în care este crescută o nouă matcă; prezența botcilor cu ouă este un semnal de pregătire a roirii. \\
Compactitatea puietului & Măsură a uniformității cu care sunt ocupate celulele cu puiet pe o ramă; un puiet compact, fără goluri, indică de regulă o matcă sănătoasă. \\
Cules & Perioada în care albinele adună nectar și polen de la plantele melifere. \\
DeepBee & Metoda din literatură pentru detecția și clasificarea automată a celulelor de fagure, integrată ca serviciu AI în BeeSmart. \\
Fagure & Structura de celule hexagonale din ceară, folosită pentru creșterea puietului și depozitarea hranei. \\
Inspecție & Operația de verificare a unui stup, în care apicultorul observă și notează starea familiei. \\
Matcă (regină) & Singura femelă fertilă a coloniei, responsabilă de ponta ouălor. \\
Nosema & Boală a albinelor adulte cauzată de paraziți interni (microsporidii); apare printre tipurile de tratament înregistrate în BeeSmart. \\
Puiet & Totalitatea stadiilor de dezvoltare a albinelor în celule: ouă, larve și puiet căpăcit. \\
Puiet căpăcit & Celulele cu larve acoperite cu un căpăcel de ceară, în stadiul de transformare în albine adulte. \\
Ramă & Cadrul de lemn pe care albinele construiesc fagurele; unitate de bază în evaluarea unui stup. \\
Roire & Procesul natural prin care o parte din colonie pleacă împreună cu matca pentru a forma o familie nouă. \\
Stup & Adăpostul unei familii de albine; în aplicație, entitatea \texttt{Hive}. \\
Stupină & Ansamblul de stupi al unui apicultor, amplasat într-o locație; în aplicație, entitatea \texttt{Apiary}. \\
Urdiniș & Deschiderea de la baza stupului prin care intră și ies albinele. \\
\textit{Varroa} & Acarianul parazit \textit{Varroa destructor}, una dintre principalele cauze ale slăbirii coloniilor; tratamentele împotriva sa sunt urmărite în aplicație. \\
\end{longtable}
